\section{Classification of the irreducible arithmetic cycles of the nonadic radical lattice}

In the arithmetic realization, a prime corresponds
to an irreducible cycle that cannot be factored into smaller positive
cycles. Nonadic periodicity may select candidates and organize
residue classes; primality remains defined by the exact multiplicative
property

\[
p=ab\Longrightarrow a=1\ \text{or}\ b=1.
\]

Once the coproduct of prime arithmetic cycles has been selected,
\cref{eq:meissel-mertens,eq:twin-prime-constant} are obtained. The helical structure and the digital root
parametrize phase and carry, but do not replace the factorization criterion.

\section{Density, regularized term, and vacancy holonomy}

The basic positive density is

\begin{equation}
\varepsilon_{\rm vac}=\log_{10}\frac{10}{9}.
\label{eq:vacancy-density}
\end{equation}

The sequence is deduced from the completion
\(1000=729+271\) determines
\begin{equation}
\varepsilon_{\rm vac}
=1-\log_{1000}(729)
=\log_{10}\frac{10}{9},
\label{eq:vacancy-density-two-forms}
\end{equation}
and the exact mechanical word
\begin{equation}
z_n=\lfloor(n+1)\varepsilon_{\rm vac}\rfloor
-\lfloor n\varepsilon_{\rm vac}\rfloor,
\qquad n\ge1.
\label{eq:vacancy-mechanical-word}
\end{equation}
Since \(\varepsilon_{\rm vac}\) is irrational, the word is Sturmian:
it has complexity \(p(N)=N+1\), balance one and topological entropy topological zero.
This dynamic differs from the full shift on the triadic
alphabet triadic: the vacancy determines an ordered boundary of zero entropy.

The suma telescopica produces
\begin{equation}
\sum_{n=1}^{N}z_n
=\lfloor(N+1)\varepsilon_{\rm vac}\rfloor,
\qquad
\sum_{n=1}^{N}(z_n-\varepsilon_{\rm vac})
=\varepsilon_{\rm vac}
-\{(N+1)\varepsilon_{\rm vac}\}.
\label{eq:vacancy-discrepancy}
\end{equation}
The partial sums of the discrepancy have modulus less than one. This
is the estimation supporting the subsequent harmonic constants.

The fraction continuous begins
\[
\varepsilon_{\rm vac}
=[0;21,1,5,1,6,2,3,\ldots].
\]
Thus the ones in \(z_n\) are separated by \(21\) or \(22\) steps.
If
\[
\sigma=22-\varepsilon_{\rm vac}^{-1}
=1-T_G(\varepsilon_{\rm vac}),
\]
then \(a_1(\sigma)=6\) and the returns of the short separation have
length \(6\) or \(7\). The pairs \(21/22\) and \(6/7\) are two scales of
the same rotation, not coincidences of integers.

For the exact word in \cref{eq:vacancy-mechanical-word}, the regularized constant is

\begin{equation}
\gamma_{\rm vac}^{+}
=\lim_{N\to\infty}\left(
\sum_{n\le N}\frac{z_n}{n}
-\varepsilon_{\rm vac}\log N
\right).
\label{eq:vacancy-gamma}
\end{equation}

The associated holonomy is

\begin{equation}
\rho_\varepsilon(k)=\exp(2\pi\ii k\varepsilon_{\rm vac}).
\label{eq:vacancy-holonomy}
\end{equation}

The three quantities have distinct types: density, finite term, and phase. The vacancy Stieltjes tower prolongs \cref{eq:vacancy-gamma}:

\begin{equation}
\gamma_{\rm vac,j}^{+}
=\varepsilon_{\rm vac}\gamma_j
+\sum_{n\ge1}\frac{(z_n-\varepsilon_{\rm vac})(\log n)^j}{n}.
\label{eq:vacancy-stieltjes}
\end{equation}

For \(j=0\), one recovers \(\gamma_{\rm vac}^{+}\). Convergence requires the
discrepancy bound for the integral development of vacancies.

Convergence is proved here. For \(j\ge0\), the sequence
\(b_n=(\log n)^j/n\) tends to zero and is decreasing from
\(n>e^j\) onward. The uniform bound \cref{eq:vacancy-discrepancy} and Dirichlet's criterion prove the convergence of
\[
\sum_{n\ge1}(z_n-\varepsilon_{\rm vac})b_n.
\]
For \(j=0\), separating
\(\varepsilon_{\rm vac}\sum_{n\le N}n^{-1}\) gives exactly
\cref{eq:vacancy-gamma}. Abel summation conservatively bounds the
tail after \(N\) by \(2/(N+1)\): one term comes from the boundary and
another from the total variation of \(1/n\). The declared exhaustive execution uses
\(N=10^8\), enumerates the Beatty positions of the ones, checks the
gap sets \(\{21,22\}\) and returns \(\{6,7\}\), and obtains as its
finite part
\[
-0.11207107505876366224207105242526397929\ldots.
\]
Combining the analytic bound with the recorded rounding and boundary
guards, the finite certificate encloses
\begin{equation}
-0.112071094143613634
<\gamma_{\rm vac}^{+}
<-0.112071055516338732.
\label{eq:vacancy-gamma-interval}
\end{equation}
This interval is a terminal value with controlled error. It does not define the
word; the word and its discrepancy precede it.

A second evaluation using multiprecision arithmetic and directed
rounding control provides, to eighteen decimal places, the certified interval
\[
-0.112071086058763562
<\gamma_{\rm vac}^{+}
<-0.112071064058763762,
\]
contained strictly within \cref{eq:vacancy-gamma-interval}. The inclusion of
both intervals verifies the stability of the finite part with respect to
truncation, the boundary guards, and the precision regime. The
result depends solely on the recurrences, the tail bounds, and the
interval arithmetic set out in this chapter.

\section{Operative content of the helicoidal representation}

The helical representation encodes a periodic phase together with a
memory increment. It has operatorial content when there exists a cocycle

\[
c(\gamma\eta)=c(\gamma)+c(\eta)
\]

that transports the carry and whose monodromy is interlinked with the
arithmetic operator. In the absence of that map, the helix constitutes only a
phase parametrization; with the interlinker, it distinguishes phase return
from state return.

\begin{remark}[Negative control]
Cofinality of cylinders does not imply dynamical
conjugacy. A counterexample to commutation between
the TPK shift and \(T_G\) precludes the intertwiner, but does not affect
\cref{eq:levy,eq:gauss-entropy}. In the vacancy sector, the divergence of
\cref{eq:vacancy-stieltjes} under the declared bound refutes the family of
regularized constants.
\end{remark}

\subsection*{Logical Scope and Domain of Validity}
Cofinality, Gauss identities, Levy, entropy, and Lochs: \emph{Exact internal theorem in the transported system}. GKW and dynamic intertwiner: \emph{Open program with falsification criterion}. Density, constant, and vacancy holonomy: \emph{Exact internal theorem with discrepancy control}.
